%=======================================================================
%  CHAPTER 3 — NUMERICAL PROBLEMS AND SOLUTIONS
%=======================================================================
\section*{\color{acc}Ch 3 \textperiodcentered{} Radio Wave Propagation \gap
\textcolor{sub}{Numerical Problems \& Solutions}}
\vspace{-4pt}{\color{acc}\rule{\textwidth}{1.1pt}}\vspace{4pt}

{\color{sub}\footnotesize \mk{This is the heaviest numerical chapter in the subject.}
A propagation numerical has appeared in \textbf{every one of the 22 papers}.
Every problem below is a real PYQ, worked in full.}

\lead{The formula sheet}
\begin{itemize}
  \item \textbf{Free space (Friis)} \;
        $P_r(d) = \dfrac{P_t G_t G_r \lambda^2}{(4\pi)^2 d^2 L}$, \qquad $\lambda = c/f$
  \item \textbf{Free space path loss} \;
        $PL = 20\log\!\left(\dfrac{4\pi d}{\lambda}\right)
            = 32.44 + 20\log f_{\text{MHz}} + 20\log d_{\text{km}}$
  \item \textbf{Power flux density} \; $P_d = \dfrac{P_t G_t}{4\pi d^2} = \dfrac{|E|^2}{120\pi}$,
        \qquad \textbf{aperture} \; $A_e = \dfrac{G_r\lambda^2}{4\pi}$, \qquad $P_r = P_d A_e$
  \item \textbf{Matched receiver} \; $V_{\text{rms}} = \sqrt{P_r R_{\text{ant}}}$, \;
        open-circuit $V_{\text{ant}} = 2V_{\text{rms}}$
  \item \textbf{Two-ray} \; $P_r = P_t G_t G_r \dfrac{h_t^2 h_r^2}{d^4}$, \;
        $PL = 40\log d - 20\log h_t - 20\log h_r$ (unity gains)
  \item \textbf{Log-distance} \;
        $PL(d) = PL(d_0) + 10\,n\log\!\left(\dfrac{d}{d_0}\right) + X_\sigma$
  \item \textbf{Okumura} \; $L_{50} = L_F + A_{mu} - G(h_{te}) - G(h_{re}) - G_{\text{AREA}}$,
        with $G(h_{te}) = 20\log\frac{h_{te}}{200}$, \;
        $G(h_{re}) = 10\log\frac{h_{re}}{3}$ for $h_{re}\le 3$ m, \;
        $20\log\frac{h_{re}}{3}$ for $3 < h_{re} < 10$ m
  \item \textbf{Hata (urban)} \;
        $L_{50} = 69.55 + 26.16\log f_c - 13.82\log h_{te} - a(h_{re})
                  + (44.9 - 6.55\log h_{te})\log d$
  \item \textbf{Delay spread} \;
        $\bar{\tau} = \dfrac{\sum P_k \tau_k}{\sum P_k}$, \;
        $\overline{\tau^2} = \dfrac{\sum P_k \tau_k^2}{\sum P_k}$, \;
        $\sigma_\tau = \sqrt{\overline{\tau^2} - \bar{\tau}^2}$
  \item \textbf{Coherence} \; $B_c \approx \dfrac{1}{50\sigma_\tau}$ (90\,\%), \;
        $B_c \approx \dfrac{1}{5\sigma_\tau}$ (50\,\%), \;
        $f_m = \dfrac{v}{\lambda}$, \; $T_c \approx \dfrac{0.423}{f_m}$
\end{itemize}

{\color{sub}\footnotesize \mk{Units discipline, the source of most lost marks:}
$P$ in W for Friis, dBm only after conversion. $d$ in \textbf{metres} in
$20\log(4\pi d/\lambda)$ but in \textbf{km} in the 32.44 form and in Hata. $f$ in
\textbf{MHz} in the 32.44 form and in Hata. Always write the unit beside every number.}

\hr

\T{Problem 1 \quad Free space: received power, E-field, flux density, rms voltage}
{\color{sub}\footnotesize \tF{3} \yr{\texttt{80 Ba}, \texttt{81 Ba} \m{8}, \textbf{78 Ch}
\m{2+2+2+2}} \gap$\cdot$\gap \mk{the classic four-part free space numerical}}

\begin{asked}{2080 Baishakh \textperiodcentered{} Q3 (b) \textperiodcentered{} [8]}
Assume free space propagation, a receiver is located 10 km away from a 50 w transmitter. The
carrier frequency is 900 MHza. antenna gain at transmitter and receiver are 1 and 2
respectively, find:
\begin{enumerate}[label=\roman*), leftmargin=1.9em]
  \item The power received at the receiver
  \item The magnitude of E-field at the receiver antenna
  \item The power flux density
  \item The rms voltage applied to the receiver input.
\end{enumerate}
\par\vspace{1pt}
The receiver antenna has a purely real impedance of 50 $\Omega$ and is matched to the receiver.
\end{asked}
\begin{asked}{2078 Chaitra \textperiodcentered{} Q3 (b) \textperiodcentered{} [2+2+2+2]}
Assume free space propagation, a receiver is located 10 km away from a 50W transmitter. The
carrier frequency is 900 MHz, antenna gain at transmitter and receiver is 1 and 2 respectively.
Find: (i) The power received at the receiver. (ii) The magnitude of the E-field at the
receiver. (iii) The rms voltage applied to the receiver input. (iv) The power flux density.
\par\vspace{1pt}
{\color{sub}Same numbers, parts (iii) and (iv) swapped. ``900 MHza'' is the 80 Ba paper's typo.}
\end{asked}

\lead{Step 0: wavelength}
\[ \lambda = \frac{c}{f} = \frac{3\times10^{8}}{900\times10^{6}} = 0.3333\ \text{m} \]

\lead{(i) Received power}
\begin{itemize}
  \item Friis, with $L = 1$ (no system loss stated):
        \[ P_r = \frac{P_t G_t G_r \lambda^2}{(4\pi)^2 d^2}
               = \frac{50 \times 1 \times 2 \times (0.3333)^2}{(4\pi)^2 (10^4)^2} \]
        \[ = \frac{50 \times 2 \times 0.11111}{157.91 \times 10^{8}}
          = \frac{11.111}{1.5791\times10^{10}} = 7.0362\times10^{-10}\ \text{W} \]
  \item In dBm:
        \[ P_r = 10\log_{10}\!\left(7.0362\times10^{-10}\times10^{3}\right)
               = \boxed{-61.53\ \text{dBm}} \; = -91.53\ \text{dBW} \]
\end{itemize}

\lead{(iii) Power flux density \mk{(do this before the E-field, it feeds it)}}
\[ P_d = \frac{P_t G_t}{4\pi d^2} = \frac{50 \times 1}{4\pi (10^4)^2}
       = \frac{50}{1.2566\times10^{9}} = \boxed{3.9789\times10^{-8}\ \text{W/m}^2} \]

\lead{(ii) Magnitude of the E-field}
\begin{itemize}
  \item Free space intrinsic impedance $\eta = 120\pi = 377\ \Omega$, and $P_d = |E|^2/\eta$:
        \[ |E| = \sqrt{P_d \times 377} = \sqrt{3.9789\times10^{-8} \times 377}
               = \sqrt{1.5\times10^{-5}} \]
        \[ = \boxed{3.873\times10^{-3}\ \text{V/m}} = 3.873\ \text{mV/m} \]
\end{itemize}

\lead{Cross-check: $P_r$ from $P_d$ and the aperture}
\[ A_e = \frac{G_r\lambda^2}{4\pi} = \frac{2 \times 0.11111}{12.566} = 0.017684\ \text{m}^2 \]
\[ P_r = P_d A_e = 3.9789\times10^{-8} \times 0.017684 = 7.036\times10^{-10}\ \text{W}
   \quad\checkmark \text{ matches (i)} \]

\lead{(iv) rms voltage at the receiver input}
\begin{itemize}
  \item The antenna is \textbf{matched} to a purely real 50 $\Omega$ load, so the power
        delivered is
        \[ P_r = \frac{V_{\text{rms}}^2}{R_{\text{ant}}}
           \;\Longrightarrow\;
           V_{\text{rms}} = \sqrt{P_r R_{\text{ant}}} \]
        \[ = \sqrt{7.0362\times10^{-10} \times 50} = \sqrt{3.5181\times10^{-8}} \]
        \[ = \boxed{1.8757\times10^{-4}\ \text{V} = 187.57\ \mu\text{V}} \]
  \item \mk{If the question asks for the open-circuit antenna voltage} $V_{\text{ant}}$ instead,
        the matched load halves it, so
        \[ V_{\text{ant}} = 2 V_{\text{rms}} = 375.13\ \mu\text{V} \]
        \mk{Read the wording: ``applied to the receiver input'' is $V_{\text{rms}}$;
        ``induced in the antenna'' is $V_{\text{ant}}$.}
\end{itemize}

\hr

\T{Problem 2 \quad Free space link with dBi antennas}
{\color{sub}\footnotesize \tP{1} \yr{\textbf{\texttt{80 Bh}} \m{6}}}

\begin{asked}{2080 Bhadra \textperiodcentered{} Q2 (a) \textperiodcentered{} [6]}
A wireless communication transmitter has an output power of 165 watts at a carrier frequency of
325 MHz. It is connected to an antenna with a gain of 12 dBi. The receiving antenna is 15 km
away and has a gain of 6 dBi. Calculate the power delivered to the receiver, considering
free-space propagation. (Assume that there are no other losses or mismatches in the system.)
\end{asked}

\begin{itemize}
  \item \mk{Gains are in dB, so work the whole problem in dB.} Converting them to linear and
        using Friis directly also works but is slower and more error-prone.
  \item \textbf{Step 1: transmit power in dBm.}
        \[ P_t = 10\log_{10}(165 \times 10^{3}) = 10\log_{10}(1.65\times10^{5})
               = 52.17\ \text{dBm} \]
  \item \textbf{Step 2: free space path loss.} Using the km / MHz form:
        \[ PL = 32.44 + 20\log_{10}(325) + 20\log_{10}(15) \]
        \[ = 32.44 + 50.24 + 23.52 = 106.20\ \text{dB} \]
        {\color{sub}\footnotesize Check with $\lambda = 3\times10^8/325\times10^6 = 0.9231$ m:
        $20\log(4\pi \times 15000/0.9231) = 106.20$ dB \checkmark}
  \item \textbf{Step 3: link equation.}
        \[ P_r = P_t + G_t + G_r - PL = 52.17 + 12 + 6 - 106.20 \]
        \[ = \boxed{-36.03\ \text{dBm}} \]
  \item \textbf{Step 4: convert back to watts if asked.}
        \[ P_r = 10^{-36.03/10}\ \text{mW} = 2.4966\times10^{-4}\ \text{mW}
               = \boxed{249.7\ \text{nW}} \]
\end{itemize}

\hr

\T{Problem 3 \quad Full link budget with feeder and mismatch loss}
{\color{sub}\footnotesize \tP{1} \yr{\textbf{72 Ash} \m{8}}}

\begin{asked}{2072 Ashwin \textperiodcentered{} Q3 \textperiodcentered{} [8]}
A BS transmitter has a power output of 10 watts operating at a frequency of 250 MHz. The
transmitter is connected by 20 m of an RF coaxial cable, which has a loss of 3-dB/100 m
specification, to an antenna that has a gain of 9 dBi. The receiving antenna is 25 km away and
has a gain of 4 dBi. There is negligible loss in the receiver feeder line, but the receiver is
mismatched; the receiving antenna and feeder cable are designed for 50 ohm impedance. The
receiver impedance loss due to mismatch is of about 0.2 dB. Calculate the power delivered to
the receiver, assuming free-space propagation.
\end{asked}

\begin{itemize}
  \item \textbf{Step 1: transmit power in dBm.}
        \[ P_t = 10\log_{10}(10 \times 10^{3}) = 40.00\ \text{dBm} \]
  \item \textbf{Step 2: transmit feeder loss.}
        \[ L_{\text{coax}} = 20\ \text{m} \times \frac{3\ \text{dB}}{100\ \text{m}}
                           = 0.60\ \text{dB} \]
  \item \textbf{Step 3: free space path loss at 25 km, 250 MHz.}
        \[ PL = 32.44 + 20\log_{10}(250) + 20\log_{10}(25) \]
        \[ = 32.44 + 47.96 + 27.96 = 108.36\ \text{dB} \]
  \item \textbf{Step 4: assemble the budget.} Every gain adds, every loss subtracts.
\end{itemize}

\vspace{2pt}
\begin{center}\small
\begin{tabular}{@{}lrl@{}}
\toprule
\hrow \thead{Item} & \thead{Value} & \thead{Note} \\
\midrule
Transmitter output      & $+40.00$ dBm & 10 W \\
Transmit feeder (coax)  & $-0.60$ dB   & 20 m at 3 dB/100 m \\
Transmit antenna gain   & $+9.00$ dB   & 9 dBi \\
Free space path loss    & $-108.36$ dB & 25 km, 250 MHz \\
Receive antenna gain    & $+4.00$ dB   & 4 dBi \\
Receive feeder loss     & $-0.00$ dB   & stated negligible \\
Receiver mismatch loss  & $-0.20$ dB   & 50 $\Omega$ system \\
\midrule
\textbf{Power at receiver} & $\mathbf{-56.16}$ \textbf{dBm} & \\
\bottomrule
\end{tabular}
\end{center}

\vspace{2pt}
\[ P_r = 40.00 - 0.60 + 9.00 - 108.36 + 4.00 - 0.20 = \boxed{-56.16\ \text{dBm}} \]
\[ = 10^{-56.16/10}\ \text{mW} = 2.42\times10^{-6}\ \text{mW} = \boxed{2.42\ \text{nW}} \]

{\color{sub}\footnotesize \mk{Marking tip:} draw the budget as a table like this one. Examiners
give method marks for a laid-out budget even when one entry is wrong.}

\hr

\T{Problem 4 \quad Radio coverage range from a receiver threshold}
{\color{sub}\footnotesize \tP{1} \yr{\textbf{70 Bh} \m{6}}}

\begin{asked}{2070 Bhadra \textperiodcentered{} Q4 \textperiodcentered{} [6]}
Determine the radio coverage range of a base station that transmits a RF signal at 150 W, given
the receiver threshold level is $-104$ dBm. Assume that the path loss at the first meter is
15 dB in a mobile radio propagation condition. (Path loss exponent $=$ 4)
\end{asked}

\begin{itemize}
  \item \textbf{Step 1: transmit power in dBm.}
        \[ P_t = 10\log_{10}(150 \times 10^{3}) = 51.76\ \text{dBm} \]
  \item \textbf{Step 2: maximum allowable path loss.}
        The link works as long as $P_r \ge -104$ dBm, so
        \[ PL_{\max} = P_t - P_{r,\min} = 51.76 - (-104) = 155.76\ \text{dB} \]
  \item \textbf{Step 3: log-distance model with $d_0 = 1$ m.}
        \[ PL(d) = PL(d_0) + 10\,n\log_{10}\!\left(\frac{d}{d_0}\right)
                 = 15 + 40\log_{10}(d) \]
  \item \textbf{Step 4: solve for $d$.}
        \[ 155.76 = 15 + 40\log_{10}(d) \]
        \[ \log_{10}(d) = \frac{155.76 - 15}{40} = \frac{140.76}{40} = 3.519 \]
        \[ d = 10^{3.519} = \boxed{3304\ \text{m} \approx 3.30\ \text{km}} \]
  \item \mk{Sanity check:} at 3.3 km with $n=4$ the loss beyond 1 m is
        $40\log(3304) = 140.8$ dB, plus the 15 dB first-metre loss gives 155.8 dB,
        which is exactly the budget. \checkmark
\end{itemize}

\hr

\T{Problem 5 \quad Monopole aperture and two-ray received power}
{\color{sub}\footnotesize \tP{1} \yr{\textbf{71 Bh} \m{6}}}

\begin{asked}{2071 Bhadra \textperiodcentered{} Q4 \textperiodcentered{} [6]}
A mobile is located 5 km away from a base station and a vertical $\lambda/4$ monopole antenna
with a gain of 2.55 dB to receive cellular radio signals. The electric field at 1 km from the
transmitter is measured to be $10^{-3}$ V/m. The carrier frequency used for this system is
900 MHz.
\begin{enumerate}[label=\alph*), leftmargin=1.6em]
  \item Find the length and effective aperture of the receiving antenna.
  \item Find the received power at the mobile using two ray ground reflection model assuming
        the height of the transmitting antenna is 50 m and the receiving antenna is 1.5 m
        above ground.
\end{enumerate}
\end{asked}

\lead{(a) Length and effective aperture}
\begin{itemize}
  \item \[ \lambda = \frac{3\times10^{8}}{900\times10^{6}} = 0.3333\ \text{m} \]
  \item \textbf{Length} of a quarter-wave monopole:
        \[ L = \frac{\lambda}{4} = \frac{0.3333}{4} = \boxed{0.0833\ \text{m} = 8.33\ \text{cm}} \]
  \item \textbf{Gain} in linear terms:
        \[ G_r = 10^{2.55/10} = 1.7989 \]
  \item \textbf{Effective aperture:}
        \[ A_e = \frac{G_r \lambda^2}{4\pi} = \frac{1.7989 \times 0.11111}{12.566}
               = \boxed{0.0159\ \text{m}^2} \]
\end{itemize}

\lead{(b) Received power from the two-ray model}
\begin{itemize}
  \item \textbf{Step 1: the two-ray E-field.} For $d$ well beyond the crossover distance,
        \[ |E_{\text{tot}}(d)| = \frac{2 E_0 d_0}{d} \cdot \frac{2\pi h_t h_r}{\lambda d} \]
        with the reference $E_0 d_0 = 10^{-3} \times 1000 = 1$ V (measured at $d_0 = 1$ km).
  \item \textbf{Step 2: substitute.}
        \[ |E_{\text{tot}}| = \frac{2 \times 1}{5000}
           \cdot \frac{2\pi \times 50 \times 1.5}{0.3333 \times 5000} \]
        \[ = 4\times10^{-4} \times \frac{471.24}{1666.7}
          = 4\times10^{-4} \times 0.28274 \]
        \[ = \boxed{1.131\times10^{-4}\ \text{V/m}} \]
  \item \textbf{Step 3: convert field to power through the aperture.}
        \[ P_r = \frac{|E_{\text{tot}}|^2}{120\pi} \cdot A_e
               = \frac{(1.131\times10^{-4})^2}{377} \times 0.0159 \]
        \[ = \frac{1.2791\times10^{-8}}{377} \times 0.0159
          = 3.3928\times10^{-11} \times 0.0159 \]
        \[ = \boxed{5.397\times10^{-13}\ \text{W} = -92.68\ \text{dBm}} \]
  \item \mk{Check the model is valid here:} crossover distance
        $d_c = 4\pi h_t h_r/\lambda = 4\pi(50)(1.5)/0.3333 = 2827$ m.
        Since $d = 5000\ \text{m} > d_c$, the $1/d^2$ field approximation holds. \checkmark
\end{itemize}

\hr

\T{Problem 6 \quad Free space path loss of a cordless phone}
{\color{sub}\footnotesize \tF{2} \yr{\textbf{77 Ch} \m{6}, \textbf{80 Ch} \m{2+4}}}

\begin{asked}{2077 Chaitra \textperiodcentered{} Q3 \textperiodcentered{} [6]}
Given a cordless phone operating at 52000 kHz frequency which has a range of 50m. Determine the
free space path loss.
\end{asked}
\begin{asked}{2080 Chaitra \textperiodcentered{} Q4 (b) \textperiodcentered{} [2+4]}
Define Doppler's effect in Small Scale Fading. (i) Given a cordless phone operating at
52000 kHz frequency which has a range of 50m. Determine the free space path loss.
\end{asked}

\begin{itemize}
  \item \textbf{Step 1: sort out the units.}
        \[ f = 52000\ \text{kHz} = 52\ \text{MHz} = 52\times10^{6}\ \text{Hz} \]
        \mk{The ``52000 kHz'' is deliberately awkward. Convert it first.}
  \item \textbf{Step 2: wavelength.}
        \[ \lambda = \frac{3\times10^{8}}{52\times10^{6}} = 5.769\ \text{m} \]
  \item \textbf{Step 3: path loss.}
        \[ PL = 20\log_{10}\!\left(\frac{4\pi d}{\lambda}\right)
              = 20\log_{10}\!\left(\frac{4\pi \times 50}{5.769}\right) \]
        \[ = 20\log_{10}(108.9) = \boxed{40.74\ \text{dB}} \]
  \item \textbf{Cross-check} with the km / MHz form, $d = 0.05$ km:
        \[ PL = 32.44 + 20\log(52) + 20\log(0.05) = 32.44 + 34.32 - 26.02 = 40.74\ \text{dB}
           \quad\checkmark \]
  \item \mk{Note how low this is} compared with the 100 dB-plus figures elsewhere: short range and
        a long wavelength both help. That is why cordless phones need only milliwatts.
\end{itemize}

\hr

\T{Problem 7 \quad dBm, dBW and received power at two distances}
{\color{sub}\footnotesize \tP{1} \yr{\texttt{81 Ba} \m{1+1+2+2}}}

\begin{asked}{2081 Baishakh \textperiodcentered{} Q3 (b) \textperiodcentered{} [1+1+2+2]}
If a transmitter produces 50 watts of power, express the transmit power in units of (i) dBm,
and (ii) dBW. If 50 watts is applied to a unity gain antenna with a 900 MHz carrier frequency,
find the received power in dBm at a free space distance of 100 m from the antenna. What is
$P_r$ (10 km)? Assume unity gain for the receiver antenna.
\end{asked}

\begin{itemize}
  \item \textbf{(i)} \[ P_t\,[\text{dBm}] = 10\log_{10}(50 \times 10^{3})
        = 10\log_{10}(5\times10^{4}) = \boxed{46.99\ \text{dBm}} \]
  \item \textbf{(ii)} \[ P_t\,[\text{dBW}] = 10\log_{10}(50) = \boxed{16.99\ \text{dBW}} \]
        {\color{sub}\footnotesize Always $\text{dBm} = \text{dBW} + 30$. \checkmark}
  \item \textbf{(iii) At $d = 100$ m}, $\lambda = 0.3333$ m, $G_t = G_r = 1$:
        \[ P_r = \frac{50 \times 1 \times 1 \times (0.3333)^2}{(4\pi)^2 (100)^2}
               = \frac{5.5556}{157.91 \times 10^{4}} = 3.5181\times10^{-6}\ \text{W} \]
        \[ = 10\log_{10}(3.5181\times10^{-3}\ \text{mW}) = \boxed{-24.54\ \text{dBm}} \]
  \item \textbf{(iv) At $d = 10$ km}, use the $1/d^2$ scaling instead of recomputing:
        \[ P_r(10\ \text{km}) = P_r(100\ \text{m}) \left(\frac{100}{10000}\right)^{2}
           = 3.5181\times10^{-6} \times 10^{-4} \]
        \[ = 3.5181\times10^{-10}\ \text{W} = \boxed{-64.54\ \text{dBm}} \]
  \item \mk{The shortcut in dB:} the distance grew by a factor 100, ie 2 decades, and free space
        costs \textbf{20 dB per decade}, so $-24.54 - 40 = -64.54$ dBm. \checkmark
\end{itemize}

\hr

\T{Problem 8 \quad Far-field (Fraunhofer) distance}
{\color{sub}\footnotesize \tP{1} \yr{\textbf{79 Ch} \m{3+3}}}

\begin{asked}{2079 Chaitra \textperiodcentered{} Q4 (b) \textperiodcentered{} [3+3]}
Explain basic propagation models. Compute the far field distance for an antenna with maximum
dimension of 1 m and operating frequency of 900 MHz.
\end{asked}

\begin{itemize}
  \item \[ \lambda = \frac{3\times10^{8}}{900\times10^{6}} = 0.3333\ \text{m} \]
  \item \[ d_f = \frac{2D^2}{\lambda} = \frac{2 \times (1)^2}{0.3333}
             = \boxed{6\ \text{m}} \]
  \item \textbf{Verify both side conditions:}
  \begin{itemize}
    \item $d_f \gg D$: \; $6 \gg 1$ \checkmark
    \item $d_f \gg \lambda$: \; $6 \gg 0.3333$ \checkmark
  \end{itemize}
  \item \mk{Meaning:} the Friis equation may only be used beyond 6 m from this antenna. Closer in,
        the wavefront is not yet planar and the formula is meaningless.
\end{itemize}

\hr

\T{Problem 9 \quad Okumura model: median path loss and received power}
{\color{sub}\footnotesize \tF{4} \yr{74 Ma \m{4+4}, \textbf{\texttt{80 Bh}} \m{4+2+2},
\textbf{\texttt{82 Bh}} \m{5}} \gap$\cdot$\gap \mk{the single most repeated numerical in the
chapter}}

\begin{asked}{2080 Bhadra \textperiodcentered{} Q3 (b) \textperiodcentered{} [4+2+2]}
Determine the mean path loss using Okumura's model for $d = 50$ km, $h_{te} = 100$ m,
$h_{re} = 10$ m in a suburban environment. If the base station transmitter radiates an EIRP of
1kW at a carrier frequency of 900 MHz. Find EIRP (dBm) and the power at the receiver where the
gain at receiving antenna is 10 dB. The medium attenuation factor is 43 dB and the gain due to
area is 9 dB for a 50 km distance and 900 MHz frequency of operation.
\end{asked}
\begin{asked}{2082 Bhadra \textperiodcentered{} Q3 (b) \textperiodcentered{} [5]}
Find the mean path loss using Okumura's Model for $d = 50$ km, $h_{te} = 100$m, $h_{re} = 10$m
in a suburban environment. If the BS transmitter radiates an EIRP oh 1 kW at a carrier frequency
of 900 MHz, find EIRP in dBm and the power at the receiver where gain of receiving antenna is
10 dB. (Given: $A_{mu}(f,d) = 43$ dB, $G_{Area}(f) = 9$ dB)
\end{asked}
\begin{asked}{2074 Magh \textperiodcentered{} Q3 \textperiodcentered{} [4+4]}
Explain Okumura model of outdoor radio propagation. Determine the median path loss for T-R
separation of 50Km, transmit antenna effective height of 100m and receive antenna effective
height of 10m in a suburban area with correction factor of 9dB at 900MHz. Assume median
attenuation relative to free space 43dB.
\end{asked}
{\color{sub}\footnotesize One set of numbers, three papers. ``EIRP oh 1 kW'' is the 82 Bh
paper's own typo.}

\lead{Step 1: free space path loss at 50 km, 900 MHz}
\[ \lambda = 0.3333\ \text{m} \]
\[ L_F = 20\log_{10}\!\left(\frac{4\pi \times 50\times10^{3}}{0.3333}\right)
       = 20\log_{10}(1.885\times10^{6}) = 125.51\ \text{dB} \]
{\color{sub}\footnotesize Check: $32.44 + 20\log(900) + 20\log(50) = 32.44 + 59.08 + 33.98
= 125.50$ dB \checkmark}

\lead{Step 2: antenna height gain factors}
\begin{itemize}
  \item Base station, $h_{te} = 100$ m (valid, since $30 < 100 < 1000$ m):
        \[ G(h_{te}) = 20\log_{10}\!\left(\frac{100}{200}\right) = 20\log_{10}(0.5)
                     = -6.02\ \text{dB} \]
        \mk{Negative, because 100 m is below the 200 m reference.}
  \item Mobile, $h_{re} = 10$ m. Since $3 < h_{re} < 10$ m the \textbf{20 log} branch applies:
        \[ G(h_{re}) = 20\log_{10}\!\left(\frac{10}{3}\right) = 20\log_{10}(3.333)
                     = +10.46\ \text{dB} \]
        \mk{Choosing the wrong branch here is the standard error. Check $h_{re}$ against 3 m
        first.}
\end{itemize}

\lead{Step 3: assemble the Okumura result}
\[ L_{50} = L_F + A_{mu} - G(h_{te}) - G(h_{re}) - G_{\text{AREA}} \]
\[ = 125.51 + 43 - (-6.02) - (10.46) - 9 \]
\[ = 125.51 + 43 + 6.02 - 10.46 - 9 = \boxed{155.07\ \text{dB}} \]

\lead{Step 4 \yr{(80 Bh, 82 Bh)}: EIRP and received power}
\begin{itemize}
  \item \[ \text{EIRP} = 1\ \text{kW} = 1000\ \text{W}
        \;\Rightarrow\; 10\log_{10}(1000 \times 10^{3}) = \boxed{60\ \text{dBm}} \]
  \item \[ P_r = \text{EIRP} - L_{50} + G_r = 60 - 155.07 + 10 = \boxed{-85.07\ \text{dBm}} \]
  \item \[ = 10^{-85.07/10}\ \text{mW} = 3.11\times10^{-9}\ \text{mW} = 3.11\ \text{pW} \]
  \item \mk{Interpretation:} $-85$ dBm is a usable but not comfortable level. Typical GSM
        sensitivity is around $-104$ dBm, so this link has about 19 dB of margin.
\end{itemize}

\hr

\T{Problem 10 \quad Okumura run backwards: find the T--R separation}
{\color{sub}\footnotesize \tP{1} \yr{\textbf{77 Ch} \m{8}}}

\begin{asked}{2077 Chaitra \textperiodcentered{} Q4 (b) \textperiodcentered{} [8]}
Employing the Okumura model compute the transmitter and receiver separation distance if median
loss is 167 dB when the carrier frequency is 2.1 GHz. Assume height of transmitting antenna is
40 m, height of receiving antenna is 2 m, for a large city.
[$A_{mu} = 34$ dB, $G_{area} = 0$ dB]
\end{asked}

\begin{itemize}
  \item \textbf{Step 1: the height gain factors.}
        \[ G(h_{te}) = 20\log_{10}\!\left(\frac{40}{200}\right) = 20\log_{10}(0.2)
                     = -13.98\ \text{dB} \]
        \[ h_{re} = 2\ \text{m} \le 3\ \text{m}
           \;\Rightarrow\; G(h_{re}) = 10\log_{10}\!\left(\frac{2}{3}\right) = -1.76\ \text{dB} \]
        \mk{Here the 10 log branch applies. Compare with Problem 9, where it was 20 log.}
  \item \textbf{Step 2: rearrange Okumura for $L_F$.}
        \[ L_{50} = L_F + A_{mu} - G(h_{te}) - G(h_{re}) - G_{\text{area}} \]
        \[ L_F = L_{50} - A_{mu} + G(h_{te}) + G(h_{re}) + G_{\text{area}} \]
        \[ = 167 - 34 + (-13.98) + (-1.76) + 0 = \boxed{117.26\ \text{dB}} \]
  \item \textbf{Step 3: invert the free space formula.}
        \[ \lambda = \frac{3\times10^{8}}{2.1\times10^{9}} = 0.1429\ \text{m} \]
        \[ L_F = 20\log_{10}\!\left(\frac{4\pi d}{\lambda}\right)
           \;\Longrightarrow\;
           d = \frac{\lambda}{4\pi}\,10^{L_F/20} \]
        \[ = \frac{0.1429}{12.566} \times 10^{117.26/20}
          = 0.011370 \times 10^{5.863} \]
        \[ = 0.011370 \times 7.2934\times10^{5} = \boxed{8292\ \text{m} \approx 8.29\ \text{km}} \]
  \item \mk{Sanity check:} substituting $d = 8.29$ km back gives
        $L_F = 32.44 + 20\log(2100) + 20\log(8.292) = 32.44 + 66.44 + 18.38 = 117.26$ dB
        \checkmark
\end{itemize}

\hr

\T{Problem 11 \quad Hata model, and comparison with free space}
{\color{sub}\footnotesize \tP{1} \yr{\texttt{79 Bh} \m{4+4}}}

\begin{asked}{2079 Bhadra \textperiodcentered{} Q3 (b) \textperiodcentered{} [4+4]}
In mobile propagation in a cellular system, find the correction factor and pathloss for a medium
size city assuming carrier frequency as 950 MHz, height of transmitting antenna at base station
is 45 m, propagation distance between antennas is 10 km and height of receiving antenna in
mobile station is 5 m. Compute free space pathloss and compare it with Hata pathloss.
\end{asked}

\begin{itemize}
  \item \textbf{Step 1: mobile antenna correction factor}, small/medium city branch:
        \[ a(h_{re}) = (1.1\log_{10} f_c - 0.7)h_{re} - (1.56\log_{10} f_c - 0.8) \]
        \[ \log_{10}(950) = 2.9777 \]
        \[ = (1.1 \times 2.9777 - 0.7)(5) - (1.56 \times 2.9777 - 0.8) \]
        \[ = (3.2755 - 0.7)(5) - (4.6452 - 0.8) \]
        \[ = 2.5755 \times 5 - 3.8452 = 12.8775 - 3.8452 = \boxed{9.032\ \text{dB}} \]
  \item \textbf{Step 2: the four Hata terms.} $\log_{10}(45) = 1.6532$, $\log_{10}(10) = 1$.
\end{itemize}

\vspace{2pt}
\begin{center}\small
\begin{tabular}{@{}llr@{}}
\toprule
\hrow \thead{Term} & \thead{Evaluation} & \thead{Value (dB)} \\
\midrule
Constant                 & $69.55$                             & $+69.550$ \\
$26.16\log f_c$          & $26.16 \times 2.9777$               & $+77.897$ \\
$-13.82\log h_{te}$      & $-13.82 \times 1.6532$              & $-22.847$ \\
$-a(h_{re})$             & $-9.032$                            & $-9.032$ \\
$(44.9-6.55\log h_{te})\log d$ & $(44.9 - 10.828)\times 1$     & $+34.072$ \\
\midrule
\textbf{$L_{50}$ (urban)} & & $\mathbf{+149.64}$ \\
\bottomrule
\end{tabular}
\end{center}

\begin{itemize}
  \item \[ \boxed{L_{50}(\text{urban}) = 149.64\ \text{dB}} \]
  \item \textbf{Step 3: free space path loss at the same distance.}
        \[ PL_{FS} = 32.44 + 20\log(950) + 20\log(10) = 32.44 + 59.56 + 20 = 112.00\ \text{dB} \]
  \item \textbf{Step 4: compare.}
        \[ L_{50} - PL_{FS} = 149.64 - 112.00 = \boxed{37.64\ \text{dB}} \]
  \item \mk{Interpretation:} the urban clutter costs almost \textbf{38 dB} over free space, ie a
        factor of about 5800 in power. This gap is exactly why free space is useless for cell
        planning and empirical models exist.
  \item \textbf{Validity check:} $f_c = 950$ MHz is inside 150--1500 MHz \checkmark,
        $d = 10$ km inside 1--20 km \checkmark, $h_{te} = 45$ m inside 30--200 m \checkmark,
        $h_{re} = 5$ m inside 1--10 m \checkmark. \mk{State this; it carries a mark.}
\end{itemize}

\hr

\T{Problem 12 \quad Hata run backwards: maximum cell radius and area}
{\color{sub}\footnotesize \tP{1} \yr{73 Ma \m{6}}}

\begin{asked}{2073 Magh \textperiodcentered{} Q3 (b) \textperiodcentered{} [6]}
Let us consider a medium sized city and assume the typical GSM downlink parameters. The Base
Station (BS) is transmitting with power 50W. The minimum acceptable received power at Mobile
Station (MS) is $-91$ dBm. The carrier frequency is 900 MHz, the height of BS is 30m and height
of MS is 1m. Estimate the maximum cell radius and corresponding cell area using Hata Model.
\par\vspace{1pt}
{\color{sub}The paper prints the full Hata model box in its header: $L_{50}$(urban)
$= 69.55 + 26.16\log f_c - 13.82\log h_{te} - a(h_{re}) + (44.9 - 6.55\log h_{te})\log d$, with
$a(h_{re}) = (1.1\log f_c - 0.7)h_{re} - (1.56\log f_c - 0.8)$ dB for a medium sized city.}
\end{asked}

\begin{itemize}
  \item \textbf{Step 1: allowable path loss.}
        \[ P_t = 10\log_{10}(50\times10^{3}) = 46.99\ \text{dBm} \]
        \[ PL_{\max} = 46.99 - (-91) = 137.99\ \text{dB} \]
  \item \textbf{Step 2: correction factor}, medium city, $\log_{10}(900) = 2.9542$:
        \[ a(h_{re}) = (1.1 \times 2.9542 - 0.7)(1) - (1.56 \times 2.9542 - 0.8) \]
        \[ = (3.2496 - 0.7) - (4.6086 - 0.8) = 2.5496 - 3.8086 = \boxed{-1.259\ \text{dB}} \]
        \mk{A negative correction is normal for a 1 m mobile antenna.}
  \item \textbf{Step 3: split Hata into $L_{50} = A + B\log d$.} With $\log_{10}(30) = 1.4771$:
        \[ A = 69.55 + 26.16(2.9542) - 13.82(1.4771) - (-1.259) \]
        \[ = 69.55 + 77.281 - 20.413 + 1.259 = 127.678 \]
        \[ B = 44.9 - 6.55(1.4771) = 44.9 - 9.675 = 35.225 \]
  \item \textbf{Step 4: solve for $d$.}
        \[ 137.99 = 127.678 + 35.225\log_{10} d \]
        \[ \log_{10} d = \frac{137.99 - 127.678}{35.225} = \frac{10.312}{35.225} = 0.29274 \]
        \[ d = 10^{0.29274} = \boxed{1.962\ \text{km}} \]
  \item \textbf{Step 5: cell area.} Treating the cell as a regular hexagon of radius $R = d$:
        \[ A_{\text{cell}} = 2.598\,R^2 = 2.598 \times (1.962)^2 = 2.598 \times 3.850
           = \boxed{10.00\ \text{km}^2} \]
        {\color{sub}\footnotesize If the cell is taken as a circle instead,
        $\pi R^2 = 12.10$ km$^2$. State which you used.}
  \item \mk{Note $d = 1.96$ km is just below the Hata validity floor of 1 km}, so the answer sits
        at the edge of the model's range. Mention this: it is worth a mark.
\end{itemize}

\hr

\T{Problem 13 \quad Hata in a large city}
{\color{sub}\footnotesize \tP{1} \yr{\textbf{73 Bh} \m{7}}}

\begin{asked}{2073 Bhadra \textperiodcentered{} Q3 (b) \textperiodcentered{} [7]}
Determine the propagation path loss for a radio signal 900 MHz cellular system operating in a
large urban city, with a base station transmitter antenna height of 100 m and mobile receiver
antenna height of 2m. The mobile unit is located at a distance of 4 km. Use the Hata propagation
path loss model.
\par\vspace{1pt}
{\color{sub}Hints printed with the question: $L_{50} = 69.55 + 26.16\log f_c - 13.82\log h_t -
\alpha(h_r) + (44.9 - 6.55\log h_t)\log d$; \ $\alpha(h_r) = (1.1\log f_c - 0.7)h_r -
(1.56\log f_c - 0.8)$ for a small to medium city, $8.29(\log 1.54h_r)^2 - 1.1$ for a large city
with $f_c \le 300$ MHz, and $3.2(\log 11.75h_r)^2 - 4.97$ for a large city with
$f_c > 300$ MHz.}
\end{asked}

\begin{itemize}
  \item \textbf{Step 1: pick the correct $a(h_{re})$ branch.} Large city, $f_c = 900 > 300$ MHz,
        so
        \[ a(h_{re}) = 3.2\left(\log_{10}(11.75\,h_{re})\right)^2 - 4.97 \]
        \[ = 3.2\left(\log_{10}(11.75 \times 2)\right)^2 - 4.97
          = 3.2\left(\log_{10}(23.5)\right)^2 - 4.97 \]
        \[ = 3.2(1.3711)^2 - 4.97 = 3.2 \times 1.8799 - 4.97 = 6.0156 - 4.97
          = \boxed{1.045\ \text{dB}} \]
  \item \textbf{Step 2: the Hata terms.} $\log_{10}(900) = 2.9542$, $\log_{10}(100) = 2$,
        $\log_{10}(4) = 0.6021$.
\end{itemize}

\vspace{2pt}
\begin{center}\small
\begin{tabular}{@{}llr@{}}
\toprule
\hrow \thead{Term} & \thead{Evaluation} & \thead{Value (dB)} \\
\midrule
Constant                 & $69.55$                         & $+69.550$ \\
$26.16\log f_c$          & $26.16 \times 2.9542$           & $+77.281$ \\
$-13.82\log h_{te}$      & $-13.82 \times 2$               & $-27.640$ \\
$-a(h_{re})$             & $-1.045$                        & $-1.045$ \\
$(44.9-6.55\log h_{te})\log d$ & $(44.9 - 13.10)\times 0.6021$ & $+19.147$ \\
\midrule
\textbf{$L_{50}$ (urban)} & & $\mathbf{+137.29}$ \\
\bottomrule
\end{tabular}
\end{center}

\[ \boxed{L_{50} = 137.29\ \text{dB}} \]

{\color{sub}\footnotesize \mk{Compare with Problem 11:} there the mobile was 5 m high in a medium
city and the loss was 149.6 dB at 10 km. Here the BS is twice as high and the distance is 2.5
times shorter, so 137.3 dB is consistent.}

\hr

\T{Problem 14 \quad Two-ray Vs free space path loss}
{\color{sub}\footnotesize \tP{1} \yr{70 Ma \m{4+4}}}

\begin{asked}{2070 Magh \textperiodcentered{} Q3 \textperiodcentered{} [4+4]}
Determine the propagation path loss for signal at 800 MHZ, with a transmitting antenna height of
30 m and a receiving antenna height of 2 m, over a distance of 10 km, using two-ray mobile
point-to-point propagation model. How is it compared with that of free-space propagation path
loss model?
\end{asked}

\begin{itemize}
  \item \textbf{Step 1: check the model is valid.}
        \[ \lambda = \frac{3\times10^{8}}{800\times10^{6}} = 0.375\ \text{m} \]
        \[ d_c = \frac{4\pi h_t h_r}{\lambda} = \frac{4\pi \times 30 \times 2}{0.375}
              = \frac{754.0}{0.375} = 2011\ \text{m} = 2.01\ \text{km} \]
        Since $d = 10\ \text{km} > d_c$, the $d^{-4}$ form applies. \checkmark
        \mk{Always do this check first: below $d_c$ the two-ray model oscillates and this formula
        is wrong.}
  \item \textbf{Step 2: two-ray path loss} (unity antenna gains):
        \[ PL = 40\log_{10} d - \left(20\log_{10} h_t + 20\log_{10} h_r\right) \]
        \[ = 40\log_{10}(10\,000) - 20\log_{10}(30) - 20\log_{10}(2) \]
        \[ = 40(4) - 29.54 - 6.02 = 160 - 35.56 = \boxed{124.44\ \text{dB}} \]
  \item \textbf{Step 3: free space path loss at the same distance.}
        \[ PL_{FS} = 32.44 + 20\log(800) + 20\log(10)
                   = 32.44 + 58.06 + 20 = \boxed{110.50\ \text{dB}} \]
  \item \textbf{Step 4: comparison.}
        \[ PL_{2\text{-ray}} - PL_{FS} = 124.44 - 110.50 = \boxed{13.94\ \text{dB} \text{ worse}} \]
  \item \mk{Why the two-ray model predicts more loss:} the ground-reflected ray arrives almost
        exactly out of phase with the direct ray and \textbf{partially cancels} it. Free space
        ignores that ray entirely, so it is over-optimistic.
  \item \mk{Note also:} the two-ray answer contains no $\lambda$ term, so it is
        \textbf{frequency independent}, while the free space answer is not.
\end{itemize}

\hr

\T{Problem 15 \quad Link budget with Okumura: reverse-link distance}
{\color{sub}\footnotesize \tP{1} \yr{\textbf{76 Bh} \m{10}}}

\begin{asked}{2076 Bhadra \textperiodcentered{} Q3 (b) \textperiodcentered{} [10]}
Estimate the appropriate distance that should be maintained for reverse link between one BTS and
mobile with appropriate link budget diagram.
\begin{enumerate}[label=\roman*), leftmargin=1.9em]
  \item Mobile is connected to antenna with 20dBi gain, with a transmitting power of 15 dBm and
        a receive sensitivity of $-75$dBm.
  \item BTS is connected to antenna with 5 dBi gain, with a transmitting power of 20dBm and a
        receive sensitivity of $-80$dBm.
  \item Cables in both systems are short, with a loss of 3-dB at each side at 900 MHz frequency
        of operation.
\end{enumerate}
\par\vspace{1pt}
Use Okamura's model for mean path loss where $G$(Area) $=$ 9dB, $A_{mu} = 43$ dB,
$h_{re} = 10$m and $h_{te} = 100$m. (Link margin (reverse link) $=$ 10 dB).
\end{asked}

\begin{itemize}
  \item \textbf{Step 1: identify the reverse link.} Reverse $=$ \textbf{mobile transmits, BTS
        receives}. So use the \textbf{mobile's} TX power and the \textbf{BTS's} sensitivity.
  \item \textbf{Step 2: maximum allowable path loss on the reverse link.}
\end{itemize}

\vspace{2pt}
\begin{center}\small
\begin{tabular}{@{}lrl@{}}
\toprule
\hrow \thead{Item} & \thead{Value} & \thead{Note} \\
\midrule
Mobile TX power        & $+15$ dBm & \\
Mobile antenna gain    & $+20$ dB  & 20 dBi \\
Mobile cable loss      & $-3$ dB   & \\
BTS antenna gain       & $+5$ dB   & 5 dBi \\
BTS cable loss         & $-3$ dB   & \\
BTS sensitivity        & $-(-80)$ dBm & the link must land at or above this \\
Link margin            & $-10$ dB  & fading reserve \\
\midrule
$\mathbf{PL_{\max}}$   & $\mathbf{104}$ \textbf{dB} & \\
\bottomrule
\end{tabular}
\end{center}

\[ PL_{\max} = 15 + 20 - 3 + 5 - 3 + 80 - 10 = \boxed{104\ \text{dB}} \]

\begin{itemize}
  \item \textbf{Step 3: convert this $L_{50}$ into a free space loss.}
        \[ G(h_{te}) = 20\log\!\left(\frac{100}{200}\right) = -6.02\ \text{dB},
           \qquad G(h_{re}) = 20\log\!\left(\frac{10}{3}\right) = +10.46\ \text{dB} \]
        \[ L_F = L_{50} - A_{mu} + G(h_{te}) + G(h_{re}) + G_{\text{AREA}} \]
        \[ = 104 - 43 + (-6.02) + 10.46 + 9 = 74.44\ \text{dB} \]
  \item \textbf{Step 4: solve for $d$.}
        \[ d = \frac{\lambda}{4\pi}\,10^{L_F/20}
             = \frac{0.3333}{12.566} \times 10^{74.44/20}
             = 0.026526 \times 5271 \]
        \[ = \boxed{139.8\ \text{m} \approx 0.14\ \text{km}} \]
  \item \textbf{Step 5: check the forward link too.} Forward $=$ BTS transmits, mobile receives:
        \[ PL_{\max} = 20 + 5 - 3 + 20 - 3 + 75 - 10 = 104\ \text{dB} \]
        \mk{Identical, so the two links are perfectly balanced} and the range is the same,
        139.8 m. Quote both and say which limits.
  \item \mk{Flag the inconsistency, it earns credit:} $A_{mu} = 43$ dB is Okumura's value at
        \emph{tens of km}, not at 140 m. Using it at this range makes the model self-inconsistent.
        The honest engineering answer is that \textbf{with the stated 43 dB clutter allowance the
        link only closes over about 140 m}, and that either the clutter figure or the antenna
        gains in the question are unrealistic.
\end{itemize}

\hr

\T{Problem 16 \quad Feasibility of a 10 km link at 2.4 GHz}
{\color{sub}\footnotesize \tP{1} \yr{\textbf{74 Bh} \m{12}} \gap$\cdot$\gap
\mk{a ``prove it does not work'' question, so the answer must end with a verdict}}

\begin{asked}{2074 Bhadra \textperiodcentered{} Q3 \textperiodcentered{} [12]}
Estimate the feasibility of a 10-km wireless link in suburban area, with one access point and
one client radio, using Okumura model for path loss. The median attenuation value is 20 dB and
gain due to environment is 13 dB. The height of access point antenna is 100 m and that of client
antennal is 10 m:
\begin{enumerate}[label=\alph*., leftmargin=1.7em]
  \item Access point is connected to antenna with 5-dBi gain, with a transmitting power of
        20-dBm and a receive sensitivity of $-80$-dBm
  \item Client is connected to antenna with 20-dBi gain, with a transmitting power of 15-dBm and
        a receive sensitivity of $-75$-dBm
  \item Cables in both systems are short, with a loss of 3-dB at each side at 2.4-GHz frequency
        of operation.
\end{enumerate}
\end{asked}

\begin{itemize}
  \item \textbf{Step 1: free space loss at 10 km, 2.4 GHz.}
        \[ \lambda = \frac{3\times10^{8}}{2.4\times10^{9}} = 0.125\ \text{m} \]
        \[ L_F = 32.44 + 20\log(2400) + 20\log(10) = 32.44 + 67.60 + 20 = 120.05\ \text{dB} \]
  \item \textbf{Step 2: Okumura median path loss.}
        \[ G(h_{te}) = 20\log(100/200) = -6.02\ \text{dB}, \qquad
           G(h_{re}) = 20\log(10/3) = +10.46\ \text{dB} \]
        \[ L_{50} = 120.05 + 20 - (-6.02) - 10.46 - 13 = \boxed{122.61\ \text{dB}} \]
  \item \textbf{Step 3: forward link, AP $\rightarrow$ client.}
        \[ P_r = 20 + 5 - 3 + 20 - 3 - 122.61 = -83.61\ \text{dBm} \]
        Client sensitivity is $-75$ dBm, so the margin is
        \[ -83.61 - (-75) = \boxed{-8.61\ \text{dB}} \quad \text{(shortfall)} \]
  \item \textbf{Step 4: reverse link, client $\rightarrow$ AP.}
        \[ P_r = 15 + 20 - 3 + 5 - 3 - 122.61 = -88.61\ \text{dBm} \]
        AP sensitivity is $-80$ dBm, so the margin is
        \[ -88.61 - (-80) = \boxed{-8.61\ \text{dB}} \quad \text{(shortfall)} \]
  \item \textbf{Step 5: verdict.}
        \mk{Both directions fall 8.6 dB short, so the 10 km link is NOT feasible} as specified,
        and it fails symmetrically, so no single end can be blamed.
  \item \textbf{Step 6: what distance \emph{would} work?} Set the margin to zero:
        \[ L_{50,\max} = 122.61 - 8.61 = 114.00\ \text{dB} \]
        \[ L_F = 114.00 - 20 + (-6.02) + 10.46 + 13 = 111.44\ \text{dB} \]
        \[ d = \frac{0.125}{12.566}\,10^{111.44/20} = 0.009947 \times 3.729\times10^{5}
             = \boxed{3.71\ \text{km}} \]
  \item \textbf{How to fix it \mk{(worth stating)}:} raise the antenna gains, raise the AP height
        (each doubling of $h_{te}$ buys 6 dB), or accept a shorter link. Raising transmit power by
        8.6 dB would also do it but may breach the licence-exempt limit at 2.4 GHz.
\end{itemize}

\hr

\T{Problem 17 \quad Power delay profile: delay spread and coherence bandwidth}
{\color{sub}\footnotesize \tF{5} \yr{70 Ma \m{4+4}, 71 Ma \m{6}, \textbf{75 Bh} \m{8},
\textbf{\texttt{81 Bh}} \m{6}} \gap$\cdot$\gap \mk{four different papers, one method}}

\lead{The method, once}
\begin{itemize}
  \item \textbf{Step 1: convert every dB power to a linear ratio}, $P_k = 10^{(\text{dB})/10}$.
        \mk{Never average in dB.}
  \item \textbf{Step 2: mean excess delay}
        $\;\bar{\tau} = \dfrac{\sum P_k\tau_k}{\sum P_k}$
  \item \textbf{Step 3: second moment}
        $\;\overline{\tau^2} = \dfrac{\sum P_k\tau_k^2}{\sum P_k}$
  \item \textbf{Step 4: rms delay spread}
        $\;\sigma_\tau = \sqrt{\overline{\tau^2} - \bar{\tau}^2}$
  \item \textbf{Step 5: coherence bandwidth}
        $\;B_c = \dfrac{1}{50\sigma_\tau}$ (90\,\%) or $\dfrac{1}{5\sigma_\tau}$ (50\,\%).
\end{itemize}

\lead{17.1 \; The 70 Ma profile \yr{(\m{4+4})}}
\sbsr{0.62}{%
\begin{asked}{2070 Magh \textperiodcentered{} Q4 \textperiodcentered{} [4+4]}
Define Doppler spread. Describe the types of small scale fading based on Doppler spread.
Calculate the mean excess delay and rms delay spread for the multipath profile given below.
Estimate the 90\% and 50\% coherence bandwidth of the channel.
\par\vspace{1pt}
{\color{sub}The profile is printed as the stem plot reproduced on the right: two components
only, $-10$ dB at $\tau = 1\ \mu$s and 0 dB at $\tau = 2\ \mu$s.}
\end{asked}
\vspace{2pt}
\begin{center}\small
\begin{tabular}{@{}rrrrr@{}}
\toprule
\hrow $\tau_k$ ($\mu$s) & dB & $P_k$ & $P_k\tau_k$ & $P_k\tau_k^2$ \\
\midrule
1 & $-10$ & 0.1 & 0.1 & 0.1 \\
2 & $0$   & 1.0 & 2.0 & 4.0 \\
\midrule
\textbf{sum} & & \textbf{1.1} & \textbf{2.1} & \textbf{4.1} \\
\bottomrule
\end{tabular}
\end{center}
\begin{itemize}
  \item \[ \bar{\tau} = \frac{2.1}{1.1} = \boxed{1.909\ \mu\text{s}} \]
  \item \[ \overline{\tau^2} = \frac{4.1}{1.1} = 3.727\ \mu\text{s}^2 \]
  \item \[ \sigma_\tau = \sqrt{3.727 - (1.909)^2} = \sqrt{3.727 - 3.645}
           = \sqrt{0.0826} = \boxed{0.287\ \mu\text{s}} \]
  \item \[ B_c(90\%) = \frac{1}{50 \times 0.287\times10^{-6}} = \boxed{69.6\ \text{kHz}} \]
  \item \[ B_c(50\%) = \frac{1}{5 \times 0.287\times10^{-6}} = \boxed{695.7\ \text{kHz}} \]
  \item \textbf{Maximum excess delay} within 10 dB of the peak: the peak is at 2 $\mu$s and the
        1 $\mu$s component is exactly 10 dB down, so
        $\tau_{10\text{dB}} = 2 - 1 = 1\ \mu$s.
\end{itemize}}{%
\figT{c3_pdp_70ma.png}
\vspace{1pt}
{\color{sub}\footnotesize The profile as printed on the 70 Ma paper.}}

\lead{17.2 \; The 75 Bh profile: greatest symbol rate without an equalizer \yr{(\m{8})}}
\begin{asked}{2075 Bhadra \textperiodcentered{} Q5 \textperiodcentered{} [8]}
Determine the smallest symbol period $T_s$, and thus the greatest symbol rate that must be sent
through RF channel with given power delay profile without using an equalizer.
\par\vspace{2pt}
\begin{tabularx}{0.60\textwidth}{@{}L{2.6cm}XXXX@{}}
\toprule
\hrow \thead{Power [dB]} & 0 & 0 & $-10$ & $-20$ \\
\thead{Delay [us]} & 0 & 50 & 75 & 100 \\
\bottomrule
\end{tabularx}
\par\vspace{1pt}
Modulation provides suitable BER performance whenever $\sigma_\tau / T_S \le 0.1$.
\end{asked}

\vspace{2pt}
\begin{center}\small
\begin{tabular}{@{}rrrrr@{}}
\toprule
\hrow $\tau_k$ ($\mu$s) & dB & $P_k$ & $P_k\tau_k$ & $P_k\tau_k^2$ \\
\midrule
0   & $0$   & 1.00 & 0.0  & 0.0 \\
50  & $0$   & 1.00 & 50.0 & 2500.0 \\
75  & $-10$ & 0.10 & 7.5  & 562.5 \\
100 & $-20$ & 0.01 & 1.0  & 100.0 \\
\midrule
\textbf{sum} & & \textbf{2.11} & \textbf{58.5} & \textbf{3162.5} \\
\bottomrule
\end{tabular}
\end{center}

\begin{itemize}
  \item \[ \bar{\tau} = \frac{58.5}{2.11} = 27.725\ \mu\text{s} \]
  \item \[ \overline{\tau^2} = \frac{3162.5}{2.11} = 1498.82\ \mu\text{s}^2 \]
  \item \[ \sigma_\tau = \sqrt{1498.82 - (27.725)^2} = \sqrt{1498.82 - 768.68}
           = \sqrt{730.14} = \boxed{27.02\ \mu\text{s}} \]
  \item \textbf{Apply the BER condition.}
        \[ \frac{\sigma_\tau}{T_s} \le 0.1
           \;\Longrightarrow\; T_s \ge \frac{\sigma_\tau}{0.1} = 10\,\sigma_\tau \]
        \[ T_s \ge 10 \times 27.02\ \mu\text{s} = \boxed{270.2\ \mu\text{s}} \]
  \item \textbf{Greatest symbol rate.}
        \[ R_s = \frac{1}{T_s} = \frac{1}{270.2\times10^{-6}} = \boxed{3.70\ \text{ksymbols/s}} \]
  \item \mk{Interpretation:} this channel is so dispersive that without an equalizer the link is
        limited to under 4 ksym/s. That single number is the practical argument for equalization
        in Ch 5.
\end{itemize}

\lead{17.3 \; The 71 Ma profile: is the channel flat for AMPS and for GSM? \yr{(\m{6})}}
\begin{asked}{2071 Magh \textperiodcentered{} Q3 (b) \textperiodcentered{} [6]}
A wireless channel is characterized by the following power-delay profile:
\par\vspace{2pt}
\begin{tabularx}{0.62\textwidth}{@{}L{2.6cm}XXXX@{}}
\toprule
\hrow \thead{Power [dB]} & 0 & $-10$ & $-20$ & $-23$ \\
\thead{Delays [ns]} & 0 & 100 & 200 & 400 \\
\bottomrule
\end{tabularx}
\par\vspace{1pt}
Determine the root mean square (rms) delay spread and the 90\% coherence bandwidth of the above
channel. Is this channel flat fading or frequency selective fading for:
\begin{enumerate}[label=\roman*), leftmargin=1.9em]
  \item An AMPS system with transmission bandwidth 30 kHz?
  \item A GSM system with transmission bandwidth 200 kHz?
\end{enumerate}
\end{asked}

\vspace{2pt}
\begin{center}\small
\begin{tabular}{@{}rrrrr@{}}
\toprule
\hrow $\tau_k$ (ns) & dB & $P_k$ & $P_k\tau_k$ & $P_k\tau_k^2$ \\
\midrule
0   & $0$   & 1.0000 & 0.00   & 0 \\
100 & $-10$ & 0.1000 & 10.00  & 1000 \\
200 & $-20$ & 0.0100 & 2.00   & 400 \\
400 & $-23$ & 0.0050 & 2.00   & 798 \\
\midrule
\textbf{sum} & & \textbf{1.1150} & \textbf{14.005} & \textbf{2201.6} \\
\bottomrule
\end{tabular}
\end{center}

\begin{itemize}
  \item \[ \bar{\tau} = \frac{14.005}{1.1150} = 12.560\ \text{ns} \]
  \item \[ \overline{\tau^2} = \frac{2201.6}{1.1150} = 1974.78\ \text{ns}^2 \]
  \item \[ \sigma_\tau = \sqrt{1974.78 - (12.560)^2} = \sqrt{1974.78 - 157.76}
           = \sqrt{1817.02} = \boxed{42.63\ \text{ns}} \]
  \item \[ B_c(90\%) = \frac{1}{50\sigma_\tau} = \frac{1}{50 \times 42.63\times10^{-9}}
           = \boxed{469.2\ \text{kHz}} \]
  \item \textbf{(i) AMPS, 30 kHz.} $30\ \text{kHz} \ll 469.2\ \text{kHz}$, so
        $B_s \ll B_c$ $\rightarrow$ $\boxed{\text{FLAT fading}}$. No equalizer needed.
  \item \textbf{(ii) GSM, 200 kHz.} $200\ \text{kHz} < 469.2\ \text{kHz}$, so strictly the
        channel is still flat, but only by a factor of \textbf{2.3}. \mk{That is not
        ``much less than'', so GSM is marginal on this channel} and a real GSM receiver would
        still run its equalizer. State it exactly this way: strictly flat, practically marginal.
  \item {\color{sub}\footnotesize If the 50\,\% definition is used instead,
        $B_c = 1/(5\sigma_\tau) = 4.69$ MHz and both systems are comfortably flat.
        \mk{Always say which definition you used.}}
\end{itemize}

\lead{17.4 \; The 81 Bh profile \yr{(\m{6})}}
\begin{asked}{2081 Bhadra \textperiodcentered{} Q5 \textperiodcentered{} [4+6]}
Compare Okumura model with Hata model. Calculate mean excess delay, rms delay spread and maximum
excess delay for the multipath profile $P_r(\tau)$ given as: $P_r(\tau) = (-20$dB, $-10$dB,
$-10$dB, 0dB) for $\tau = (0, 1, 2, 5)$ second. Also estimate 50\% coherence bandwidth.
\par\vspace{1pt}
{\color{sub}\mk{``second'' is the paper's own slip}: the delays are microseconds. Work in
$\mu$s and say so.}
\end{asked}

\vspace{2pt}
\begin{center}\small
\begin{tabular}{@{}rrrrr@{}}
\toprule
\hrow $\tau_k$ & dB & $P_k$ & $P_k\tau_k$ & $P_k\tau_k^2$ \\
\midrule
0 & $-20$ & 0.01 & 0.00 & 0.00 \\
1 & $-10$ & 0.10 & 0.10 & 0.10 \\
2 & $-10$ & 0.10 & 0.20 & 0.40 \\
5 & $0$   & 1.00 & 5.00 & 25.00 \\
\midrule
\textbf{sum} & & \textbf{1.21} & \textbf{5.30} & \textbf{25.50} \\
\bottomrule
\end{tabular}
\end{center}

\begin{itemize}
  \item \[ \bar{\tau} = \frac{5.30}{1.21} = \boxed{4.380\ \text{s}} \]
  \item \[ \overline{\tau^2} = \frac{25.50}{1.21} = 21.074\ \text{s}^2 \]
  \item \[ \sigma_\tau = \sqrt{21.074 - (4.380)^2} = \sqrt{21.074 - 19.183}
           = \sqrt{1.891} = \boxed{1.375\ \text{s}} \]
  \item \textbf{Maximum excess delay (10 dB).} The peak is 0 dB at $\tau = 5$. Components within
        10 dB of it are those at $\ge -10$ dB, ie $\tau = 1, 2, 5$. So
        \[ \tau_{10\text{dB}} = 5 - 1 = \boxed{4\ \text{s}} \]
  \item \[ B_c(50\%) = \frac{1}{5 \times 1.375} = \boxed{0.1455\ \text{Hz}} \]
  \item \mk{Source defect, state it:} the paper prints the delays in \textbf{seconds}. A real
        multipath channel has delays of \textbf{microseconds}, and a 0.15 Hz coherence bandwidth
        is physically impossible. \textbf{Solve it as printed}, then add one line noting that
        with $\mu$s the answers become
        $\bar{\tau} = 4.38\ \mu$s, $\sigma_\tau = 1.375\ \mu$s, $B_c(50\%) = 145.5$ kHz.
\end{itemize}

\hr

\T{Problem 18 \quad Doppler shift and coherence time}
{\color{sub}\footnotesize \tF{4} \yr{\textbf{78 Ch}, \texttt{81 Ba}, \textbf{80 Ch},
\textbf{\texttt{82 Bh}}}}

\begin{asked}[PROBLEM]{not a PYQ \textperiodcentered{} the standard Rappaport example \added}
A US digital cellular system operating at $f_c = 900$ MHz serves a vehicle moving at 70 km/h.
Compute the received carrier frequency when the mobile moves (a)~directly towards the
transmitter, (b)~directly away from it, and (c)~perpendicular to the direction of arrival.
Hence find the Doppler spread and the 50\% coherence time.
\end{asked}
{\color{sub}\footnotesize Four papers ask for the Doppler \emph{derivation} in prose
(\yr{78 Ch}, \yr{81 Ba}, \yr{80 Ch}, \yr{82 Bh}) and none prints numbers. This worked
example is what turns that derivation into a full answer. \added}

\begin{itemize}
  \item \textbf{Step 1: convert the velocity.} \mk{Always to m/s.}
        \[ v = 70\ \frac{\text{km}}{\text{h}} \times \frac{1000}{3600}
             = 19.444\ \text{m/s} \]
  \item \textbf{Step 2: wavelength and maximum Doppler shift.}
        \[ \lambda = \frac{3\times10^{8}}{900\times10^{6}} = 0.3333\ \text{m} \]
        \[ f_m = \frac{v}{\lambda} = \frac{19.444}{0.3333} = \boxed{58.33\ \text{Hz}} \]
  \item \textbf{(a) Towards the transmitter}, $\theta = 0^\circ$, $\cos\theta = +1$:
        \[ f_d = +58.33\ \text{Hz}
           \;\Rightarrow\; f = 900\ \text{MHz} + 58.33\ \text{Hz}
           = \boxed{900.00005833\ \text{MHz}} \]
  \item \textbf{(b) Away}, $\theta = 180^\circ$, $\cos\theta = -1$:
        \[ f_d = -58.33\ \text{Hz}
           \;\Rightarrow\; f = \boxed{899.99994167\ \text{MHz}} \]
  \item \textbf{(c) Perpendicular}, $\theta = 90^\circ$, $\cos\theta = 0$:
        \[ f_d = 0 \;\Rightarrow\; f = \boxed{900\ \text{MHz exactly}} \]
  \item \textbf{Step 3: coherence time.}
        \[ T_c \approx \frac{0.423}{f_m} = \frac{0.423}{58.33} = 7.25\ \text{ms} \]
  \item \mk{Interpretation:} the channel stays essentially constant for only 7.25 ms. A GSM burst
        is 0.577 ms, so the channel is safely static within one burst, which is exactly the design
        intent. A data frame lasting longer than 7 ms would see the channel change under it.
\end{itemize}
